\documentclass{tietreport}

% Import images
\usepackage{graphicx}

% Bibliography configuration
\usepackage[
  date=year,
  style=alphabetic,
  backend=biber,
  sorting=nty,
  maxnames=5,
  minnames=3,
  maxalphanames=4,
  minalphanames=3,
  backref=true,
  doi=false,
  isbn=false,
  url=false,
  eprint=false
]{biblatex}
\addbibresource{references.bib}

%% ==================================================
%% CUSTOMIZE THESE FIELDS FOR YOUR PROJECT
%% ==================================================

% Course/Lab header
\TitlePageHeader{%
  UCS503P - Software Engineering Project%
}

% Main project title
\ProjectTitle{%
  FinShield AI%
}

% Subtitle or project description
\ProjectSubTitle{%
  Intelligent Financial Crime Detection \& Investigation Platform%
}

% Student information (up to 4 students)
\author{%
  \texttt{1024170265} Harshit Goyal \\
  \texttt{1024170057} Arnav Gupta \\
  \texttt{1024170274} Dikshit Gupta \\
  \texttt{1024170269} Jasmandeep Sharma%
}

% Group number, degree, year, program
\TitlePageSubText{%
  Group: \texttt{1} , BE Third Year, CSE%
}

% Faculty advisor/guide name
\AdvisorName{%
   Nisha Thakur%
}

% Evaluation period and department info
\TitlePageFooterText{{%
    \large Prototype Stage Evaluation \\[0.2cm]
    \bfseries Computer Science and Engineering
    Department%
  } \\[0.15cm]
  Thapar Institute of Engineering and Technology,
  Patiala%
}

%% ==================================================
\let\cleardoublepage\clearpage

\begin{document}
% Generate title page
\maketitle
\tableofcontents

% ===== CHAPTER 1: INTRODUCTION =====

\chapter{Introduction}
\section{Background and Context}

Financial fraud and money laundering involve large numbers of transactions
that may contain complex and changing patterns. Traditional approaches that
depend only on fixed rules may not be sufficient to identify all suspicious
activities. FinShield AI is proposed as an intelligent financial crime
detection and investigation platform that combines machine learning,
rule-based analysis, user behavior analysis, risk scoring, and investigation
support.

The system is designed to receive transaction data, validate and preprocess
it, analyze the transaction using multiple detection approaches, and assign
a risk level. Suspicious transactions can then be converted into alerts and
investigation cases for analysts. Explainable AI techniques and an AI-based
investigation assistant are planned to support the investigation process.

\subsection{Problem Statement}

Financial institutions must quickly detect suspicious transactions without overwhelming investigators with false positives, a common pitfall of single-method systems. FinShield AI tackles this by unifying machine learning, AML rules, behavioral tracking, and hybrid risk scoring into one platform. Beyond detection, the system delivers clear risk explanations and investigative tools to help analysts act decisively.

\section{Project Objectives}

\begin{enumerate}
  \item \textbf{Objective 1:} Develop a secure platform for managing users,
    roles, and financial transaction data.
  \item \textbf{Objective 2:} Detect suspicious financial transactions using
    machine learning, AML rules, and user behavior analysis.
  \item \textbf{Objective 3:} Combine different detection results to generate
    a meaningful risk score and identify high-risk transactions.
  \item \textbf{Objective 4:} Support fraud investigation through alerts,
    explainable predictions, case management, and investigation assistance.
  \item \textbf{Objective 5:} Provide dashboards and reports for analysts,
    compliance officers, and system administrators.
\end{enumerate}

% ===== CHAPTER 2: METHODOLOGY & DESIGN =====

\chapter{Methodology and Design}

\section{System Architecture}

FinShield AI follows a modular architecture in which the frontend interacts
with the backend through APIs. The backend manages authentication,
transactions, fraud analysis, alerts, investigations, and communication with
the database. The fraud analysis layer contains machine learning, AML rule
processing, and user behavior analysis. Their outputs are intended to be
combined by the hybrid risk scoring component.

\subsection{High-Level Design}

The high-level design separates the system into user interaction,
application services, fraud intelligence, and data storage components.
Customers/users submit transactions through the interface. The backend
validates the transaction and passes the required information to the
fraud-analysis components. The resulting risk information is stored and can
be used by analysts and compliance officers for further action.


% --------------------------------------------------
% Diagram placeholders
% --------------------------------------------------

\begin{figure}[htbp]
  \centering
  \includegraphics[width=0.88\textwidth]{gantt_chart.jpeg}
  \caption{Three-Month Project Gantt Chart}
  \label{fig:gantt}
\end{figure}

\begin{figure}[htbp]
  \centering
  \includegraphics[width=0.88\textwidth]{ucd.jpeg}
  \caption{Use Case Diagram of FinShield AI}
  \label{fig:usecase}
\end{figure}

\begin{figure}[htbp]
  \centering
  \includegraphics[width=0.88\textwidth]{activity_diagram.jpeg}
  \caption{Activity Diagram of FinShield AI}
  \label{fig:activity}
\end{figure}

\begin{figure}[htbp]
  \centering
  \includegraphics[width=0.88\textwidth]{use_case_template_1.png}
    \caption{Use Case Template: User Authentication \& Profile Management}  \label{fig:usecase_1}
\end{figure}

\begin{figure}[htbp]
  \centering
  \includegraphics[width=0.88\textwidth]{use_case_template_2.png}
  \caption{Use Case Template: Transaction Management}
  \label{fig:usecase_2}
\end{figure}

\begin{figure}[htbp]
  \centering
  \includegraphics[width=0.88\textwidth]{use_case_template_3.png}
  \caption{Use Case Template: ML Fraud Detection \& Risk Scoring}
  \label{fig:usecase_3}
\end{figure}

\begin{figure}[htbp]
  \centering
  \includegraphics[width=0.88\textwidth]{use_case_template_4.png}
  \caption{Use Case Template: Fraud Alert \& Notification}
  \label{fig:usecase_4}
\end{figure}


\begin{figure}[htbp]
  \centering
  \noindent\makebox[\textwidth][c]{%
    \includegraphics[width=0.88\textwidth]{use_case_template_5.png}%
  }
  \caption{Use Case Template: Fraud Case Investigation}
  \label{fig:usecase_5}
\end{figure}

\begin{figure}[htbp]
  \centering
  \includegraphics[width=0.88\textwidth]{use_case_template_6.png}
  \caption{Use Case Template: Explainable AI Analysis}
  \label{fig:usecase_6}
\end{figure}


\begin{figure}[htbp]
  \centering
  \includegraphics[width=0.88\textwidth]{use_case_template_7.png}
  \caption{Use Case Template: System \& Rule Management}
  \label{fig:usecase_7}
\end{figure}

\begin{figure}[htbp]
  \centering
  \includegraphics[width=0.88\textwidth]{use_case_template_8.png}
  \caption{Use Case Template: Risk Reports \& Compliance Analytics}
  \label{fig:usecase_8}
\end{figure}

\begin{figure}[htbp]
\centering
\includegraphics[width=0.88\textwidth]{collaboration_diagram.jpeg}
\caption{Collaboration Diagram of FinShield AI}
\label{fig:collaboration}
\end{figure}

\begin{figure}[htbp]
\centering
\includegraphics[width=0.88\textwidth]{sequence_diagram.jpeg}
\caption{Sequence Diagram of FinShield AI}
\label{fig:sequence}
\end{figure}

\begin{figure}[htbp]
\centering
\includegraphics[width=0.88\textwidth]{Level0_DFD.drawio.png}
\caption{Level 0 Data Flow Diagram}
\label{fig:dfd_0}
\end{figure}

\begin{figure}[htbp]
\centering
\includegraphics[width=0.88\textwidth]{Level1_DFD.drawio.png}
\caption{Level 1 Data Flow Diagram}
\label{fig:dfd_1}
\end{figure}

\begin{figure}[htbp]
\centering
\includegraphics[width=0.88\textwidth]{Level2_DFD.drawio.png}
\caption{Level 2 Data Flow Diagram}
\label{fig:dfd_2}
\end{figure}

\begin{figure}[htbp]
\centering
\includegraphics[width=0.88\textwidth]{Class_Diagram.drawio.png}
\caption{Class Diagram}
\label{fig:class}
\end{figure}

\subsection{Technical Stack}

\begin{itemize}
  \item \textbf{Framework/Language:} Python for machine learning and
    backend development, with a web-based frontend.
  \item \textbf{Database:} PostgreSQL for structured storage of users,
    transactions, alerts, investigations, and related records.
  \item \textbf{Tools:} Git and GitHub for version control and collaborative
    development; Jupyter Notebook for data analysis and model experiments.
  \item \textbf{Libraries:} Pandas and NumPy for data processing, with
    appropriate machine learning libraries for model development.
\end{itemize}

\section{Implementation Details}

The prototype implementation is being developed in modules so that each
member can work on a specific part while following the same project
architecture and common transaction data schema. The initial implementation
focuses on the system foundation and the core fraud-detection workflow.
The complete investigation, reporting, explainability, and AI-assistant
features are planned for later stages of the three-month development
timeline.

\subsection{Prototype Modules Implemented}

\subsubsection*{Authentication and RBAC}

The Authentication and Role-Based Access Control (RBAC) module provides
secure access to the platform. Users are authenticated before accessing
protected functionality, and permissions are controlled according to their
assigned roles. The planned roles include Customer/User, Analyst/Investigator,
Compliance Officer, and System Administrator. This separation ensures that
each actor can access only the functions relevant to their responsibilities.

\begin{figure}[htbp]
\centering
\includegraphics[width=0.88\textwidth]{authentication.jpeg}
\caption{Authentication and RBAC}
\label{fig:authentication_rbac}
\end{figure}


\subsubsection*{Transaction Management}

The Transaction Management module handles the submission and storage of
transaction information. A transaction is first received by the system,
validated, and prepared for subsequent fraud analysis. The module provides
the basic connection between user transaction input and the fraud-detection
pipeline.

\begin{figure}[htbp]
\centering
\includegraphics[width=0.88\textwidth]{transaction.jpeg}
\caption{Transaction Management}
\label{fig:transaction}
\end{figure}


\subsubsection*{ML Fraud Detection}

The ML Fraud Detection module uses transaction attributes to identify
potentially fraudulent transactions. The initial development includes
dataset analysis, common schema definition, preprocessing, and baseline
model development. Model performance will be evaluated using fraud-detection
metrics such as precision, recall, F1-score, and ROC-AUC rather than relying
only on accuracy.

\begin{figure}[htbp]
\centering
\includegraphics[width=0.88\textwidth]{ml.jpeg}
\caption{ML Fraud Detection}
\label{fig:ml}
\end{figure}


\subsubsection*{AML Detection}

The AML Detection module provides rule-based analysis for identifying
transactions that may indicate suspicious financial activity. The module
is designed to work with the common transaction schema so that its output
can later be combined with the machine learning and behavior-analysis
results. The initial stage focuses on defining and testing suitable AML
rules; the complete rule engine and hybrid risk-scoring workflow will be
developed in later phases.

\begin{figure}[htbp]
\centering
\includegraphics[width=0.88\textwidth]{aml.jpeg}
\caption{AML Detection}
\label{fig:aml}
\end{figure}


\subsection{Initial Data and Schema Preparation}

The project uses a common transaction-data schema across all modules.
The selected datasets are being studied and prepared so that the ML,
AML, backend, and future behavior-analysis components use a consistent
representation of transaction information. A data dictionary is maintained
to document the meaning and type of each field. Preprocessing is performed
before the data is supplied to the machine learning pipeline.

The schema is treated as a shared interface between team members. Any
required change is reviewed by the team before it is introduced, reducing
the risk of incompatibility between the ML, AML, backend, and frontend
components.

\subsection{Prototype Development Status}

At the prototype stage, the project has moved from requirement analysis and
system design toward initial implementation. The current work establishes
the foundation for the remaining development phases. The complete system
will later integrate behavior analysis, hybrid risk scoring, explainable
AI, fraud investigation, dashboards, notifications, compliance reporting,
and the AI investigation assistant.

% ===== END AT SECTION 2.2.1 AS REQUESTED =====

\end{document}
